% Spanish for Beginners — Week 1 Study Guide (Unit 1 start)
% Compile: pdflatex unit-1-intro.tex

\documentclass[10pt,a4paper]{article}
\usepackage[utf8]{inputenc}
\usepackage[T1]{fontenc}
\usepackage[margin=1.2cm]{geometry}
\usepackage{enumitem}
\setlist{nosep,leftmargin=*}
\usepackage{booktabs}
\usepackage{xcolor}
\usepackage{fancyhdr}
\usepackage{microtype}
\definecolor{w1blue}{RGB}{0,80,120}
\pagestyle{fancy}
\fancyhf{}
\fancyfoot[L]{\small Spanish for Beginners | Week 1}
\fancyfoot[R]{\small\thepage}
\raggedright

\begin{document}
\begin{center}
{\large\bfseries\color{w1blue} Week 1 -- Unit 1 Start}\\[0.2em]
\footnotesize Preguntas, ser/tener/llamarse, n\'umeros 0--15.
\end{center}\vspace{0.5em}

\section*{\color{w1blue}Preguntas de identificaci\'on}
\begin{itemize}
\item \textbf{¿C\'omo te llamas?} $\to$ Me llamo... / Francisco.
\item \textbf{¿Cu\'antos a\~nos tienes?} $\to$ Treinta y cinco.
\item \textbf{¿Cu\'al es tu apellido?} $\to$ Mart\'inez.
\item \textbf{¿De d\'onde eres?} $\to$ De Valencia.
\item \textbf{¿En qu\'e trabajas?} $\to$ Soy camarero.
\end{itemize}

\section*{\color{w1blue}Ser, tener, llamarse}
\textbf{Ser} (identity, origin, profession): soy, eres, es, somos, sois, son. \\
\textbf{Tener} (age: tengo 20 a\~nos): tengo, tienes, tiene, tenemos, ten\'eis, tienen. \\
\textbf{Llamarse} (name): me llamo, te llamas, se llama...

\section*{\color{w1blue}N\'umeros 0--15}
cero, uno, dos, tres, cuatro, cinco, \textbf{seis}, siete, ocho, nueve, diez, once, doce, trece, catorce, quince. \\
\textit{Note:} 6 = seis (not ``sies'').
\end{document}
